% sec9_2.tex — Section 9.2: Tools for Unit-Root Econometrics

The basic tool in unit-root econometrics is what is called the ``functional central
limit theorem'' (FCLT). For the FCLT to be applicable, we need to specialize I(1)
processes by placing restrictions on the associated I(0) processes. Which restrictions to place, that is, which class of I(0) processes to focus on, differs from author
to author, depending on the version of the FCLT in use. In this book, following the
practice in the literature, we focus on linear I(0) processes. After defining linear
I(0) processes, this section introduces some concepts and results that will be used
repeatedly in the rest of this chapter.

\subsection*{Linear I(0) Processes}

Our definition of linear I(0) processes is as follows.

\begin{definition}[Linear I(0) processes]
\label{def:9.2}
A linear I(0) process can be written as a
constant plus a zero-mean linear process $\{u_t\}$ such that
\begin{align}
u_t &= \psi(L)\varepsilon_t, \quad
\psi(L) \equiv \psi_0 + \psi_1 L + \psi_2 L^2 + \cdots
\quad \text{for } t = 0, \pm 1, \pm 2, \ldots.
\label{eq:9.2.1} \\
\{\varepsilon_t\} &\text{ is independent white noise (i.i.d.\ with mean 0 and }
\E(\varepsilon_t^2) \equiv \sigma^2 > 0\text{),}
\label{eq:9.2.2}
\end{align}
\begin{align}
\sum_{j=0}^{\infty} j|\psi_j| &< \infty,
\label{eq:9.2.3a}
\tag{9.2.3a} \\
\psi(1) &\neq 0.
\label{eq:9.2.3b}
\tag{9.2.3b}
\end{align}
\end{definition}

The innovation process $\{\varepsilon_t\}$ is assumed to be white noise i.i.d. This requirement is
for simplifying the exposition and can be relaxed substantially.\footnote{For
example, the innovation process can be stationary martingale differences. See,
e.g., Theorem~3.8 of Tanaka (1996, p.~80).} In the rest of this
chapter, we will refer to linear I(0) processes as simply I(0) processes, without the
qualifier ``linear.'' Also, unless otherwise stated, $\{\varepsilon_t\}$ will be an independent white
noise process and $\{u_t\}$ a zero-mean I(0) process.

Condition~\eqref{eq:9.2.3a} (sometimes called the \textbf{one-summability} condition) is
stronger than the more familiar condition of absolute summability. This stronger
assumption will make it easier to prove large sample results of the next section,
with the help of the Beveridge-Nelson decomposition to be introduced shortly.
Since $\{\psi_j\}$ is absolutely summable, the linear process $\{u_t\}$ is strictly stationary
and ergodic by Proposition~6.1(d). The $j$-th order autocovariance of $\{u_t\}$ will be
denoted $\gamma_j$. Since $\E(u_t) = 0$, $\gamma_j = \E(u_t u_{t-j})$. By Proposition~6.1(c) the autocovariances are absolutely summable.

To understand condition~\eqref{eq:9.2.3b}, recall that the long-run variance (which we
denote $\lambda^2$ in this chapter) equals the value of the autocovariance-generating function at $z = 1$ by Proposition~6.8(b). By Proposition~6.6 it is given by
\begin{equation}
\label{eq:9.2.4}
\lambda^2 \equiv \text{long-run variance of } \{u_t\}
= \gamma_0 + 2 \sum_{j=1}^{\infty} \gamma_j
= \sigma^2 \cdot [\psi(1)]^2.
\end{equation}
The condition~\eqref{eq:9.2.3b} thus ensures the long-run variance to be positive.

\subsection*{Approximating I(1) by a Random Walk}

Let $\{\xi_t\}$ be I(1) so that $\Delta \xi_t = \delta + u_t$ where $u_t \equiv \psi(L)\varepsilon_t$ is a zero-mean I(0)
process satisfying \eqref{eq:9.2.1}--\eqref{eq:9.2.3b} with $\E(\xi_0^2) < \infty$. Using the following identity
(to be verified in Review Question~1):
\begin{align}
\psi(L) &= \psi(1) + \Delta \alpha(L), \quad \Delta \equiv 1 - L, \notag \\
\alpha(L) &\equiv \sum_{j=0}^{\infty} \alpha_j L^j, \quad
\alpha_j = -(\psi_{j+1} + \psi_{j+2} + \cdots) \quad (j = 0, 1, 2, \ldots),
\label{eq:9.2.5}
\end{align}
we can write $u_t$ as
\begin{equation}
\label{eq:9.2.6}
u_t \equiv \psi(L)\varepsilon_t = \psi(1) \cdot \varepsilon_t + \eta_t - \eta_{t-1}
\quad \text{with } \eta_t \equiv \alpha(L)\varepsilon_t.
\end{equation}
It can be shown (see Analytical Exercise~7) that $\alpha(L)$ is absolutely summable.
So, by Proposition~6.1(a), $\{\eta_t\}$ is a well-defined zero-mean covariance-stationary
process (it is actually ergodic stationary by Proposition~6.1(d)). Substituting~\eqref{eq:9.2.6}
into~\eqref{eq:9.1.4}, we obtain (what is known in econometrics as) the \textbf{Beveridge-Nelson
decomposition}:
\begin{align}
\xi_t &= \delta \cdot t + \sum_{s=1}^{t}[\psi(1) \cdot \varepsilon_s + \eta_s - \eta_{s-1}] + \xi_0 \notag \\
&= \delta \cdot t + \psi(1) \sum_{s=1}^{t} \varepsilon_s + \eta_t + (\xi_0 - \eta_0)
\quad \Bigl(\text{since } \sum_{s=1}^{t}(\eta_s - \eta_{s-1}) = \eta_t - \eta_0\Bigr).
\label{eq:9.2.7}
\end{align}
Thus, any linear I(1) process can be written as the sum of a linear time trend ($\delta \cdot t$),
a driftless random walk or a stochastic trend ($\psi(1)(\varepsilon_1 + \varepsilon_2 + \cdots + \varepsilon_t)$), a stationary
process ($\eta_t$), and an initial condition ($\xi_0 - \eta_0$). Stated somewhat differently, we
have

\begin{proposition}[Beveridge-Nelson decomposition]
\label{prop:9.1}
Let $\{u_t\}$ be a zero-mean
I(0) process satisfying \eqref{eq:9.2.1}--\eqref{eq:9.2.3b}. Then
\[
u_1 + u_2 + \cdots + u_t = \psi(1)(\varepsilon_1 + \varepsilon_2 + \cdots + \varepsilon_t) + \eta_t - \eta_0,
\]
where $\eta_t \equiv \alpha(L)\varepsilon_t$, $\alpha_j = -(\psi_{j+1} + \psi_{j+2} + \cdots)$. $\{\eta_t\}$ is zero-mean ergodic
stationary.\footnote{For a more general condition, which does not require the I(0)
process $\{u_t\}$ to be linear, under which $u_1 + \cdots + u_t$ is decomposed as shown here,
see Theorem~5.4 of Hall and Heyde (1980). In that theorem, $\{\varepsilon_t\}$ is a
stationary martingale difference sequence, so the stochastic trend is a martingale.}
\end{proposition}

For a moment set $\delta = 0$ in~\eqref{eq:9.2.7} so that $\{\xi_t\}$ is driftless I(1). An important
implication of the decomposition is that any driftless I(1) process is dominated by
the stochastic trend $\psi(1) \sum_{s=1}^{t} \varepsilon_s$ in the following sense. Divide both sides of the
decomposition~\eqref{eq:9.2.7} (with $\delta = 0$) by $\sqrt{t}$ to obtain
\begin{equation}
\label{eq:9.2.8}
\frac{\xi_t}{\sqrt{t}} = \psi(1) \cdot \frac{1}{\sqrt{t}} \sum_{s=1}^{t} \varepsilon_s
+ \left[\frac{\xi_0}{\sqrt{t}} + \frac{\eta_t}{\sqrt{t}} - \frac{\eta_0}{\sqrt{t}}\right].
\end{equation}
Since $\E(\xi_0^2) < \infty$ by assumption, $\E[(\xi_0/\sqrt{t})^2] \to 0$ as $t \to \infty$. So $\xi_0/\sqrt{t}$ converges in mean square (and hence in probability) to zero. The same applies to
$\eta_t/\sqrt{t}$ and $\eta_0/\sqrt{t}$. So the terms in the brackets on the right-hand side of~\eqref{eq:9.2.8} can
be ignored asymptotically. In contrast, the first term has a limiting distribution of
$N(0, \sigma^2 \cdot [\psi(1)]^2)$ by the Lindeberg-Levy CLT of Section~2.1. In this sense, the stochastic trend grows at rate $\sqrt{t}$. Using~\eqref{eq:9.2.4}, the stochastic trend can be written as
\begin{equation}
\label{eq:9.2.9}
\psi(1)(\varepsilon_1 + \varepsilon_2 + \cdots + \varepsilon_t)
= \lambda \cdot \left(\frac{\varepsilon_1}{\sigma} + \frac{\varepsilon_2}{\sigma} + \cdots + \frac{\varepsilon_t}{\sigma}\right),
\end{equation}
which shows that changes in the stochastic trend in $\xi_t$ have a variance of $\lambda^2$, the
long-run variance of $\{\Delta \xi_t\}$.

Now suppose $\delta \neq 0$ so that $\{\xi_t\}$ is I(1) with drift. As is clear from dividing
both sides of~\eqref{eq:9.2.7} by $t$ rather than $\sqrt{t}$ and applying the same argument just given
for the $\delta = 0$ case, the stochastic trend as well as the stationary component can be
ignored asymptotically, with $\xi_t/t$ converging in probability to $\delta$. In this sense the
time trend dominates the I(1) process in large samples.

\subsection*{Relation to ARMA Models}

In Section~6.2, we defined a zero-mean stationary ARMA process $\{u_t\}$ as the
unique covariance-stationary process satisfying the stationary ARMA($p$, $q$)
equation
\begin{equation}
\label{eq:9.2.10}
\phi(L) u_t = \theta(L) \varepsilon_t,
\end{equation}
where $\phi(L)$ satisfies the stationarity condition. If $\{\varepsilon_t\}$ is independent white noise
and if $\theta(1) \neq 0$, then the ARMA($p$, $q$) process is zero-mean I(0). To see this, note
from Proposition~6.5(a) that $u_t$ can be written as~\eqref{eq:9.2.1} with
\begin{equation}
\label{eq:9.2.11}
\psi(L) = \phi(L)^{-1} \theta(L).
\end{equation}
Since $\phi(1) \neq 0$ by the stationarity condition and $\theta(1) \neq 0$ by assumption, condition~\eqref{eq:9.2.3b} (that $\psi(1) \neq 0$) is satisfied. By Proposition~6.4(a), the coefficient
sequence $\{\psi_j\}$ is bounded in absolute value by a geometrically declining sequence,
which ensures that~\eqref{eq:9.2.3a} is satisfied.

Given the zero-mean stationary ARMA($p$, $q$) process $\{u_t\}$ just described, the
associated I(1) process $\{\xi_t\}$ defined by the relation $\Delta \xi_t = u_t$ is called an \textbf{autoregressive integrated moving average} (\textbf{ARIMA($\bm{p}$, 1, $\bm{q}$)}) process. Substituting the
relation $u_t = (1 - L)\xi_t$ into~\eqref{eq:9.2.10}, we obtain
\begin{equation}
\label{eq:9.2.12}
\phi^*(L) \xi_t = \theta(L) \varepsilon_t,
\end{equation}
with $\phi^*(L) \equiv \phi(L)(1 - L)$, $\phi(L)$ stationary, and $\theta(1) \neq 0$. One of the roots of
the associated autoregressive polynomial equation $\phi^*(z) = 0$ is unity and all other
roots lie outside the unit circle. More generally, a class of I($d$) processes can be
represented as satisfying~\eqref{eq:9.2.12} where $\phi^*(L)$ is now defined as
\begin{equation}
\label{eq:9.2.13}
\phi^*(L) = \phi(L)(1 - L)^d.
\end{equation}
So $\phi^*(z) = 0$ now has a root of unity with a multiplicity of $d$ (i.e., $d$ unit roots) and
all other roots lie outside the unit circle. This class of I($d$) processes is called autoregressive integrated moving average (\textbf{ARIMA($\bm{p}$, $\bm{d}$, $\bm{q}$)}) processes. The first
parameter ($p$) refers to the order of autoregressive lags (not counting the unit
roots), the second parameter ($d$) refers to the order of integration, and the third
parameter ($q$) is the number of moving average lags. Since $\phi(L)$ satisfies the
stationarity condition, taking $d$-th differences of an ARIMA($p$, $d$, $q$) produces a
zero-mean stationary ARMA($p$, $q$) process. So an ARIMA($p$, $d$, $q$) process is
I($d$).

\subsection*{The Wiener Process}

The next two sections will present a variety of unit-root tests. The limiting distributions of their test statistics will be written in terms of \textbf{Wiener processes} (also
called \textbf{Brownian motion processes}). Some of you may already be familiar with
this from continuous-time finance, but to refresh your memory,

\begin{definition}[Standard Wiener processes]
\label{def:9.3}
A standard Wiener (Brownian
motion) process $W(\cdot)$ is a continuous-time stochastic process, associating each
date $t \in [0, 1]$ with the scalar random variable $W(t)$, such that
\begin{enumerate}
\item $W(0) = 0$;
\item for any dates $0 \leq t_1 < t_2 < \cdots < t_k \leq 1$, the changes
\[
W(t_2) - W(t_1),\; W(t_3) - W(t_2),\; \ldots,\; W(t_k) - W(t_{k-1})
\]
are independent multivariate normal with $W(s) - W(t) \sim N(0, (s-t))$ (so in
particular $W(1) \sim N(0,1)$);
\item for any realization, $W(t)$ is continuous in $t$ with probability~1.
\end{enumerate}
\end{definition}

Roughly speaking, the \textbf{functional central limit theorem} (\textbf{FCLT}), also called the
\textbf{invariance principle}, states that a Wiener process is a continuous-time analogue
of a driftless random walk. Imagine that you generate a realization of a standard
random walk (whose changes have a unit variance) of length $T$, scale the graph
vertically by deflating all its values by $\sqrt{T}$, and then horizontally compress this
normalized graph so that it fits the unit interval $[0, 1]$. In Panel~(a) of Figure~9.2,
that is done for a sample size of $T = 10$. The sample size is increased to 100 in
Panel~(b), and then to 1,000 in Panel~(c). As the figure shows, the graph becomes
increasingly dense over the unit interval. The FCLT assures us that there exists a
well-defined limit as $T \to \infty$ and that the limiting process is the Wiener process.
Property~(2) of Definition~9.3 is a mathematical formulation of the property that the
sequence of instantaneous changes of a Wiener process is i.i.d. The continuous-time analogue of a driftless random walk whose changes have a variance of $\sigma^2$,
rather than unity, can be written as $\sigma W(r)$.

\begin{figure}[htbp]
\centering
\framebox[\textwidth]{\rule{0pt}{4in}\parbox{0.8\textwidth}{\centering
  Figure 9.2: Illustration of the Functional Central Limit Theorem\\
  (a) $T = 10$ \quad (b) $T = 100$ \quad (c) $T = 1000$}}
\caption{Illustration of the Functional Central Limit Theorem}
\label{fig:9.2}
\end{figure}

To pursue the analogy between a driftless random walk and a Wiener process
a bit further, consider a ``demeaned standard random walk'' which is constructed
from a standard driftless random walk by subtracting the sample mean. That is, let
$\{\xi_t\}$ be a driftless random walk with $\Var(\Delta \xi_t) = 1$ and define
\begin{equation}
\label{eq:9.2.14}
\xi_t^\mu \equiv \xi_t - \frac{\xi_0 + \xi_1 + \cdots + \xi_{T-1}}{T}
\quad (t = 0, 1, \ldots, T-1).
\end{equation}
The continuous-time analogue of this demeaned random walk is a \textbf{demeaned standard Wiener process} defined as
\begin{equation}
\label{eq:9.2.15}
W^\mu(r) \equiv W(r) - \int_0^1 W(s)\,\mathrm{d}s.
\end{equation}
(We defined the demeaned series for $t = 0, 1, \ldots, T-1$, to match the dating
convention of Proposition~9.2 below. If the demeaned series is defined for $t =
1, 2, \ldots, T$, then its continuous-time analogue, too, is~\eqref{eq:9.2.15}.)

We can also create from the standard random walk $\{\xi_t\}$ a detrended series:
\begin{equation}
\label{eq:9.2.16}
\xi_t^\tau \equiv \xi_t - \hat{\alpha} - \hat{\delta} \cdot t
\quad (t = 0, 1, \ldots, T-1),
\end{equation}
where $\hat{\alpha}$ and $\hat{\delta}$ are the OLS estimates of the intercept and the time coefficient in the
regression of $\xi_t$ on $(1, t)$ $(t = 0, 1, \ldots, T-1)$. The continuous-time analogue of
the detrended random walk is a \textbf{detrended standard Wiener process} defined as
\begin{align}
W^\tau(r) &\equiv W(r) - a - d \cdot r, \notag \\
a &\equiv \int_0^1 (4 - 6s) W(s)\,\mathrm{d}s, \notag \\
d &\equiv \int_0^1 (-6 + 12s) W(s)\,\mathrm{d}s.
\label{eq:9.2.17}
\end{align}
The coefficients $a$ and $d$ are the limiting random variables of $\hat{\alpha}$ and $\hat{\delta}$, respectively,
as $T \to \infty$ (see, e.g., Phillips and Durlauf, 1986, for a derivation). Therefore,
if a linear time trend is fitted to a \emph{driftless} random walk by OLS, the estimated
linear time trend $\hat{\delta}$ converges in distribution to a random variable $d$; even in large
samples $\hat{\delta}$ will never be zero unless by accident. This phenomenon is sometimes
called \textbf{spurious detrending}.

\subsection*{A Useful Lemma}

Unit-root tests utilize statistics involving I(0) and I(1) processes. The key results
collected in the following proposition will be used to derive the limiting distributions of the test statistics of the unit-root tests of Sections~9.3--9.4.

\begin{proposition}[Limiting distributions of statistics involving I(0) and I(1)
variables]
\label{prop:9.2}
Let $\{\xi_t\}$ be driftless I(1) so that $\Delta \xi_t$ is zero-mean I(0) satisfying
\eqref{eq:9.2.1}--\eqref{eq:9.2.3b} and $\E(\xi_0^2) < \infty$. Let
\[
\lambda^2 \equiv \textit{long-run variance of } \{\Delta \xi_t\}
\quad \text{and} \quad
\gamma_0 \equiv \Var(\Delta \xi_t).
\]

\medskip\noindent
Then
\begin{enumerate}[label=(\alph*)]
\item $\displaystyle \frac{1}{T^2} \sum_{t=1}^{T} (\xi_{t-1})^2
\;\dto\; \lambda^2 \cdot \int_0^1 W(r)^2\,\mathrm{d}r,$

\item $\displaystyle \frac{1}{T} \sum_{t=1}^{T} \Delta\xi_t\, \xi_{t-1}
\;\dto\; \frac{\lambda^2}{2}\, W(1)^2 - \frac{\gamma_0}{2}.$
\end{enumerate}

\medskip\noindent
Let $\{\xi_t^\mu\}$ be the demeaned series created from $\{\xi_t\}$ by the formula~\eqref{eq:9.2.14}. Then
\begin{enumerate}[label=(\alph*),start=3]
\item $\displaystyle \frac{1}{T^2} \sum_{t=1}^{T} (\xi_{t-1}^\mu)^2
\;\dto\; \lambda^2 \cdot \int_0^1 [W^\mu(r)]^2\,\mathrm{d}r,$

\item $\displaystyle \frac{1}{T} \sum_{t=1}^{T} \Delta\xi_t\, \xi_{t-1}^\mu
\;\dto\; \frac{\lambda^2}{2}\big\{[W^\mu(1)]^2 - [W^\mu(0)]^2\big\} - \frac{\gamma_0}{2}.$
\end{enumerate}

\medskip\noindent
Let $\{\xi_t^\tau\}$ be the detrended series created from $\{\xi_t\}$ by the formula~\eqref{eq:9.2.16}. Then
\begin{enumerate}[label=(\alph*),start=5]
\item $\displaystyle \frac{1}{T^2} \sum_{t=1}^{T} (\xi_{t-1}^\tau)^2
\;\dto\; \lambda^2 \cdot \int_0^1 [W^\tau(r)]^2\,\mathrm{d}r,$

\item $\displaystyle \frac{1}{T} \sum_{t=1}^{T} \Delta\xi_t\, \xi_{t-1}^\tau
\;\dto\; \frac{\lambda^2}{2}\big\{[W^\tau(1)]^2 - [W^\tau(0)]^2\big\} - \frac{\gamma_0}{2}.$
\end{enumerate}

\medskip\noindent
The convergence is joint. That is, a vector consisting of the statistics indicated in
(a)--(f) converges to a random vector whose elements are the corresponding random
variables also indicated in (a)--(f).
\end{proposition}

Part~(a), for example, reads as follows: the sequence of random variables $\{(1/T)^2
\sum_{t=1}^T (\xi_{t-1})^2\}$ indexed by $T$ converges in distribution to a random variable $\lambda^2 \int W^2\,
\mathrm{d}r$. The limiting random variables in the proposition are all written in terms of
standard Wiener processes. Note that the \emph{same} Wiener process $W(\cdot)$ appears in
(a) and~(b) and that the demeaned and detrended Wiener processes in (c) and~(f) are
derived from the Wiener process appearing in (a) and~(b). So the limiting random
variables in (a)--(f) can be correlated.

To gain some understanding of this result, suppose temporarily that $\{\xi_t\}$ is
a driftless random walk with $\Var(\Delta \xi_t) = \sigma^2$. Given that a Wiener process is
the continuous-time analogue, it is not surprising that ``$\sum_{t=1}^T (\xi_{t-1})^2$'' suitably normalized by a power of $T$ becomes ``$\sigma^2 \int W^2$.'' Perhaps what you might not have
guessed is the normalization by $T^2$. One way to see why the square of $T$ provides
the right normalization, is to recall that $\E[(\xi_t)^2] = \Var(\xi_t) = \sigma^2 \cdot t$. Thus,
\begin{equation}
\label{eq:9.2.18}
\E\!\left[\sum_{t=1}^{T} (\xi_{t-1})^2\right]
= \sigma^2 \sum_{t=1}^{T} (t-1)
= \sigma^2 \cdot (T-1) \frac{T}{2}.
\end{equation}
So the mean of $\sum_{t=1}^T (\xi_{t-1})^2$ grows at rate $T^2$. In order to construct a random
variable that could have a limiting distribution, $\sum_{t=1}^T (\xi_{t-1})^2$ will have to be divided
by $T^2$.

Now suppose that $\{\xi_t\}$ is a general driftless I(1) process whose first difference
is a zero-mean I(0) process satisfying \eqref{eq:9.2.1}--\eqref{eq:9.2.3b}. Then serial correlation in
$\Delta \xi_t$ can be taken into account by replacing $\sigma^2$ by $\lambda^2$ (the long-run variance of
$\{\Delta \xi_t\}$). This is due to implication~\eqref{eq:9.2.9} of the Beveridge-Nelson decomposition that a driftless I(1) process is dominated in large samples by a random walk
whose changes have a variance of $\lambda^2$. Put differently, the limiting distribution of
$\sum_{t=1}^T (\xi_{t-1})^2 / T^2$ would be the same if $\{\xi_t\}$, instead of being a general driftless
I(1) process with serially correlated first differences, were a driftless random walk
whose changes have a variance of $\lambda^2$.

For a proof of Proposition~9.2, see, e.g., Stock (1994, pp.~2751--2753). It is a
rather mechanical application of the FCLT and a theorem called the ``continuous
mapping theorem.'' Since $W(1) \sim N(0, 1)$, the limiting random variable in (b) is
\begin{equation}
\label{eq:9.2.19}
\left(\frac{\lambda^2}{2}\right) X - \frac{\gamma_0}{2}
\end{equation}
where $X \sim \chi^2(1)$ (chi-squared with 1 degree of freedom). Proof of (b) can be done
without the fancy apparatus of the FCLT and the continuous mapping theorem and
is left as Analytical Exercise~1.

\bigskip
\noindent\textsc{questions for review} \hrulefill

\begin{enumerate}
\item (Verifying Beveridge-Nelson for simple cases) \enspace Verify~\eqref{eq:9.2.5} for $\psi(L) = 1 +
\psi_1 L$. Do the same for $\psi(L) = (1 - \phi L)^{-1}$ where $|\phi| < 1$.

\item (A CLT for I(0)) \enspace Proposition~6.9 is about the asymptotics of the sample mean
of a linear process. Verify that the paragraph right below Proposition~9.1 is a
proof of Proposition~6.9 with the added assumption of one-summability.

\item (Is the stationary component in Beveridge-Nelson I(0)?) \enspace Consider a zero-mean I(0) process
\[
u_t = \psi_0 \varepsilon_t - 2\varepsilon_{t-1} + \varepsilon_{t-2}.
\]
Verify that the long-run variance of $\{\eta_t\}$ defined in~\eqref{eq:9.2.6} is zero.

\item (Initial values and demeaned values) \enspace Let $\{\xi_t\}$ be driftless I(1) so that $\xi_t$ can
be written as $\xi_t = \xi_0 + u_1 + u_2 + \cdots + u_t$. Does the value of $\xi_0$ affect the
value of $\xi_t^\mu$ where $\{\xi_t^\mu\}$ is the demeaned series created from $\{\xi_t\}$ by the formula~\eqref{eq:9.2.14}? [Answer: No.]

\item (Linear trends and detrended values) \enspace Let $\xi_t$ be I(1) with drift, so it can be
written as the sum of a linear time trend $\alpha + \delta \cdot t$ and a driftless I(1) process.
Do the values of $\alpha$ and $\delta$ affect the value of $\xi_t^\tau$ where $\{\xi_t^\tau\}$ is the detrended
series created from $\{\xi_t\}$ by the formula~\eqref{eq:9.2.16}? [Answer: No.]
\end{enumerate}
